\section{Verified transformation geometry as a proposal-only search projection}
\label{sec:vtg}

The assurance architecture above determines when a mathematical claim may change durable authority. It intentionally does not identify proof search with theorem truth. This section sharpens that separation for a new, still-unresolved research question: whether verifier-defined mathematical reachability admits a useful \emph{local} geometry at all. The formulation below distinguishes exact operational semantics, partial knowledge of that semantics, path equivalence, representation change, intrinsic path cost, budget-conditioned control, and the cost of constructing the geometry itself. Those distinctions are required before geometric language can carry mathematical meaning.

\subsection{Frozen subject and operational-state sufficiency}
Fix a formal target $P$, assumptions $\Gamma$, a registered logic/checker $L$, an exact environment/version $E$, an operator basis $\Omega$, an active representation $R$, a root quantity of interest $Q$, and a registered admissibility/cost profile $A$. Write the frozen solver subject as
\[
\Theta=(P,\Gamma,L,E,\Omega,R,Q,A).
\]
Any change to a load-bearing coordinate defines a new subject unless an explicit transport certificate licenses reuse.

A state geometry is well defined only if the operational state records all future-relevant information. Let $h_1$ and $h_2$ be two execution histories that are encoded as the same state $x$. The encoding is \emph{operationally sufficient} when, under frozen $\Theta$, the legal successor relation and transition-cost law from $x$ are independent of whether $x$ was reached through $h_1$ or $h_2$. If proof-state metadata, mutable environment, caches, accumulated assumptions, solver memory or another history-dependent coordinate changes future legal moves, that coordinate must either be included in the operational state or retained in an explicitly separate control/resource state. Without this condition, a purported quantity $d(x,y)$ is history dependent and therefore not a function of states.

\subsection{Exact transition semantics, discovered map and solution region}
Let
\[
\Sigma_{\Theta}(x,T,y)\in\{\PASS,\FAIL,\CC\}
\]
be the exact transition semantics: whether applying the typed transformation $T$ to state $x$ yields $y$ under the registered verifier. This semantic relation is distinct from the finite materialized subcomplex $\mathcal C^+_{\Theta,t}$ containing only states and transitions actually discovered and checked by time $t$.

For an interactive theorem prover, two assurance levels must not be conflated. A \emph{kernel-derivation edge} carries a local proof object whose mathematical derivation is independently checkable. A \emph{replay-validated operational edge} records an exact tactic-state transition under a bound proof-state snapshot, tool/version subject and replay receipt. The latter is useful for search reproducibility but remains operational rather than theorem authority; a proposed edge has neither status until the relevant check passes. Thus a tactic-state transition can be a verified road in the search map without being treated as an independently authoritative theorem step. Final mathematical authority remains attached to the assembled proof object and its ordinary audited kernel receipt.

The operational map is therefore a partial-belief object
\[
\mathcal B_{\Theta,t}=(K^+_t,K^-_t,K^?_t,U_t,\Gamma_{\mathrm{cov},t}),
\]
where $K^+$ are materialized assurance-bound transitions, $K^-$ are scoped refutations or blockers, $K^?$ are candidate/unverified transitions, $U$ represents explicitly unknown/unexplored regions and $\Gamma_{\mathrm{cov}}$ records the coverage contract. The binding rule is
\[
\boxed{\text{unknown transition}\;\neq\;\text{blocked transition}.}
\]
Failure to find a route in an incomplete $\mathcal B_{\Theta,t}$ is therefore a search result, not a mathematical impossibility result.

Even a certified complete \emph{registered operational map} supports only a scoped statement: no route exists under the bound subject, chart and operator basis. It does not by itself establish theorem falsity, mathematical non-existence, or even proof-system unprovability. Lifting such a result to unprovability would require a separate completeness theorem connecting the registered operator basis to the proof calculus; theorem refutation requires ordinary mathematical evidence such as a proof of negation or countermodel.

The verified solution region is defined independently of any geometry:
\[
\mathcal V_P=\{x:\operatorname{OriginalSemanticsCheck}(x,P,\Gamma,L)=\PASS\}.
\]
No distance, embedding, field, search frequency or model confidence can put a state into $\mathcal V_P$.

\subsection{Verified paths, higher-order proof structure and congruent path quotienting}
Let $\operatorname{Path}(\Theta)$ denote the free category generated by materialized verified elementary transitions: states are objects, verified finite paths are morphisms, empty paths are identities and path concatenation is composition. This microscopic object is directed and symbolic; no useful geometry is assumed.

Proof discovery and proof certification remain distinct. A solver may follow a chronological trajectory
\[
s_0\to s_1\to\cdots\to s_n,
\]
while the final proof certificate is a dependency DAG, proof term or typed hypergraph. \texttt{TaskEpisode} preserves chronology; the proof DAG/certificate preserves the mathematical dependency object. A solution-assembly receipt is proposal/assurance plumbing and cannot itself grant theorem authority.

State canonicalization also leaves path redundancy. A path equivalence relation $\sim$ is safe for quotienting only when it is an endpoint-preserving congruence:
\[
\gamma\sim\delta\Longrightarrow s(\gamma)=s(\delta),\quad t(\gamma)=t(\delta),
\]
and, whenever the compositions are defined,
\[
\gamma\sim\delta\Longrightarrow \alpha\gamma\beta\sim\alpha\delta\beta.
\]
Thus equivalence is not merely similarity of traces. For concurrency reduction, Orion permits an adjacent exchange
\[
uabv\sim ubav
\]
only when a scoped verifier-bound witness establishes that $a$ and $b$ are independent/commuting transformations in the relevant context. Absence of a declared dependency is not evidence of independence. More general proof-path identifications may later be generated by certified joinability or proof-term-equivalence diagrams; the Phase-0 implementation does not assume that every such higher cell has already been discovered. The quotient $\operatorname{Path}(\Theta)/{\sim}$ therefore removes only witnessed redundancy.

Path equivalence is semantic, not automatically economic: equivalent paths can have different cache behavior, wall time or verification cost. A search quotient must retain the relevant cost frontier unless cost invariance is separately established.

\subsection{Representations, portals, exact quotients and abstractions}
The solver representation system is more general than a differential-geometric atlas. Let $\mathsf{Rep}_P$ be a transformation network/category whose objects are problem representations and whose arrows are certified transformations such as exact equivalences, one-way reductions, quotients, bounded approximations, relaxations and dualities. The exact invertible arrows form the atlas/groupoid-like portion; irreversible reductions should not be called coordinate changes without qualification.

A target-conditioned semantic quotient can preserve the root QoI while still changing reachability. Therefore a quotient used for navigation requires a separate \emph{navigation quotient contract}. For a quotient $q:X\to\bar X$, an exact reachability-preserving claim requires at least target-label preservation and suitable two-way transition conditions. An intentionally over-approximating abstraction is a different object: it may expose an abstract route that has no concrete realization. Such a route must be concretized/replayed before use or returned as a spurious-route residual that triggers refinement. A forward simulation may therefore justify a sound over-approximation without justifying exact reachability equivalence. If the quotient or abstraction is used to estimate route cost, the distortion relation between abstract and concrete costs must also be stated or left unidentified.

\subsection{Hard admissibility and intrinsic path-cost geometry}
A mathematical route first passes non-compensatory admissibility. Let
\[
\Lambda(e)\in\{\PASS,\FAIL,\CC\}
\]
encode hard requirements such as specification alignment, licensed assumptions, valid representation portals, trusted verifier semantics and root-scope preservation. Cheap computation cannot compensate for $\Lambda(e)=\FAIL$.

For admissible transitions, path cost lives in a registered algebra $A$ rather than an assumed universal scalar. Different coordinates may compose differently: compute/time may be additive, peak memory may use a maximum, assumptions naturally compose by set union, trust may use a worst-case/meet operation, and evidence carries provenance rather than a generic numeric sum. When the chosen carrier and composition support an enriched/generalized metric construction, define the intrinsic operator/representation-relative path value
\[
D_{\Theta}(x,y)=\inf_{\gamma:x\rightsquigarrow y} C_A(\gamma).
\]
If the task does not justify one scalar/ordered algebra, Orion retains a Pareto or set-valued reachability frontier and does not promote it to a scalar distance by convention.

Crucially, the resource budget is not part of the intrinsic geometric axioms. Budget $B$ defines feasible reachability/control:
\[
x\Rightarrow_B y
\quad\Longleftrightarrow\quad
\exists\gamma:x\rightsquigarrow y\text{ with }C_A(\gamma)\preceq B.
\]
A budget-truncated distance can violate triangle inequality: two six-unit legs can each be feasible under a ten-unit cutoff while their twelve-unit composition is mapped to infinity. Goal-relative $V_{P,B}$ is therefore a control/value object above the intrinsic path geometry, not automatically a metric.

\subsection{Non-vacuity, constructibility and geometry certification}
Bare existence of a greedily navigable geometry is not a useful scientific claim. In any finite connected unweighted graph with target $t$, the exact graph distance
\[
h_t(v)=d_G(v,t)
\]
has the property that every non-target vertex has a neighbor $u$ on a shortest path with $h_t(u)=h_t(v)-1$. Thus global shortest-path computation can always be hidden in a perfect local navigation function.

\begin{proposition}[Oracle-distance trivialization]
For a finite connected graph and fixed target, exact target distance induces a well-founded strict local descent to the target. Consequently, an existence-only claim for local navigability is non-informative unless the cost and information used to construct/query the geometry are constrained.
\end{proposition}
\begin{proof}
Choose any shortest path from a non-target state $v$ to $t$. Its first edge reaches a neighbor whose distance to $t$ is exactly one smaller. Because graph distance is a nonnegative integer, repeated strict descent terminates at distance zero, namely at $t$.
\end{proof}

A VTG geometry must therefore carry a \emph{nontriviality/constructibility contract}. At minimum it records construction compute and verifier cost, coordinate or label size per state, global auxiliary storage, local distance/action query cost, target dependence, whether gold routes or gold distances were used, expected reuse, invalidation/staleness, route/search stretch and held-out transfer. We distinguish, for example, exact finite enumeration, polynomial or fixed-parameter construction on a registered class, amortized precomputation, approximate learned construction, evaluator-only oracle geometry and unknown/general intractability. An evaluator-only gold-distance geometry remains a valuable upper bound but is not eligible as evidence for a nontrivial solver improvement.

A candidate routing geometry is a derived, version-bound artifact
\[
\mathcal G_{\Theta,t}=\operatorname{Geom}(\mathcal B_{\Theta,t}),
\]
bound to the specification/QoI, verifier/environment subject, operator basis, representation, map revision, cost algebra and construction version. It must also declare the strongest certification class actually established, for example:
\[
\begin{gathered}
\texttt{EXACT\_COST\_TO\_GO},\quad
\texttt{ADMISSIBLE\_LOWER\_BOUND},\quad
\texttt{CONSISTENT\_HEURISTIC},\\
\texttt{EMPIRICAL\_RANKER},\quad
\texttt{UNCERTIFIED}.
\end{gathered}
\]
An empirical ranker can prioritize search but cannot inherit optimality/completeness claims that require an admissible or consistent heuristic.

\subsection{Certified navigation basins and learned-geometry support}
A useful geometry need not be globally navigable. Let $S$ be a registered \emph{navigation basin}: a subset on which the representation, legal-action relation and progress certificate are known to apply. The strongest local-navigation statement can be made theorem-like without assuming Euclidean geometry. Let $W$ be well founded and let $h:S\to W$. A local policy $\pi$ is \emph{descent-complete} on $S$ when every non-goal $x\in S$ has a verified legal successor $\pi(x)\in S$ with
\[
h(\pi(x))\prec h(x),
\]
and every minimal reachable $h$-state in $S$ is a verified goal. Then repeated application of $\pi$ reaches a goal in finitely many steps, because an infinite strictly descending chain in $W$ is impossible. The basin contract must state its boundary/escape behavior; it is not a claim of global navigability. For real-valued scores on infinite spaces, strict descent alone is not sufficient; a well-founded rank, uniform/proper progress condition or another termination certificate is required.

This separates three progressively stronger notions:
\begin{enumerate}
    \item \emph{local congruence}: lower score tends to rank useful legal successors higher;
    \item \emph{closed-loop navigability}: a local/bounded-local policy reaches verified goals on the registered class;
    \item \emph{certified descent}: a well-founded progress argument proves termination for the registered basin.
\end{enumerate}
The Phase-0/1 programme initially tests the first two; it does not presume the third exists for general mathematics.

A learned geometry also requires an explicit support receipt. Training subjects, behavior/search policy, state/edge sampling process, label source, coverage coordinates, unseen operators/charts/scales, structural assumptions, hidden-route leakage checks, fresh-test subjects, OOD diagnostics and staleness diagnostics are recorded separately. This prevents a self-confirming loop in which the current policy avoids a region, collects no data there, and then treats absence of data as evidence that the region is unpromising. A learned geometry can support a fresh scoped empirical claim only relative to the registered support and leakage contract; it cannot self-certify exact global reachability.

\subsection{Monotonicity, quantifiers, composition and authority noninterference}
Two simple monotonicity facts delimit what a map failure can mean. Under a frozen subject with unchanged edge semantics/costs, expanding the materialized verified map cannot remove an already verified route or worsen the best known route frontier. Likewise, if $\Omega\subseteq\Omega'$ and old operators retain their semantics, the reachable set under $\Omega'$ contains the reachable set under $\Omega$. Consequently, no-route under the current operator basis cannot imply no-route under a future enlarged basis.

In deterministic Phase 0, the primitive claim is existential path reachability. If the framework later admits stochastic, nondeterministic or adversarial actions, every reachability statement must state its quantifier semantics: existence of one path, controllability by a policy, robustness for all admissible outcomes, almost-sure reachability, a probability lower bound, an expected-cost bound, or a game/adversary guarantee. These claims are not interchangeable and require their corresponding transition-kernel/control semantics.

Local proof objects also require global amalgamation before parent closure. Locally verified children do not imply a verified parent merely because each child passes independently. A composition receipt must bind child certificate hashes, overlap/substitution consistency, discharged assumptions, parent invariants and the final parent verifier. This is the formal counterpart of GLUE: local sections are candidates for composition, not votes for global truth.

Let $\pi_{\mathrm{math}}$ denote mathematical authority and let $G$ be any update that changes only map estimates, quotient views, geometry, field values, search conductance or route priority. The implementation contract remains
\[
\pi_{\mathrm{math}}(G(x))=\pi_{\mathrm{math}}(x).
\]
Thus a monotonically decreasing energy, short geometric route or highly reinforced proof path remains proposal-level until the exact proof/specification/trust gates pass.

\begin{proposition}[Map, geometry and quotient noninterference]
Suppose (i) a route is called verified only when every materialized transition on it has the registered operational or local-derivation assurance, (ii) incomplete map coverage cannot produce an impossibility or unprovability conclusion, (iii) path/representation quotients grant no authority and any abstract route that requires lifting is revalidated in the original semantics, and (iv) geometry/dynamics have no write access to mathematical authority. Then map incompleteness, erroneous geometry, or a spurious abstract route can increase search cost or cause search failure, but cannot by itself promote or refute the target theorem.
\end{proposition}
\begin{proof}
By (i), any route admitted as operationally verified has concrete assurance-bound transitions rather than an untyped prediction. By (ii), missing roads cannot be converted into a negative theorem statement. By (iii), quotient membership or abstract reachability is not accepted as a concrete proof: a spurious route fails lifting/revalidation. Finally, (iv) blocks any direct promotion from routing state. Hence all three failure modes remain inside the search/control projection until an ordinary mathematical assurance gate is satisfied.
\end{proof}

\subsection{Fieldability is closed-loop, not local correlation}
A useful geometry must be evaluated as a policy substrate, not by the correlation of one-step scores alone. The registered \texttt{FieldabilityProfile} therefore separates construction cost, reuse break-even, invalidation/staleness, local action alignment, strict-greedy completion, bounded-branch completion, route stretch, false-descent/trap rate, verifier calls and out-of-distribution failure. The nontriviality contract adds the representation size, auxiliary storage and local query costs needed to rule out a hidden routing table.

\begin{figure}[t]
\centering
\includegraphics[width=0.49\linewidth]{../../../paper/figures/generated/unified_local_vs_closed_loop.pdf}\hfill
\includegraphics[width=0.49\linewidth]{../../../paper/figures/generated/unified_path_quotient.pdf}
\caption{Development known-world stresses for two different search claims. \textbf{Left:} a constructed heuristic is locally correct on most decision points yet strict greedy descent terminates in a dead end; retaining a bounded frontier finds the valid route. \textbf{Right:} $m$ pairwise-independent transformations have $m!$ sequential histories but one partial-order computation class \emph{when independence is actually certified}. These figures test failure semantics and search redundancy; they are not evidence that a learned theorem-proving geometry already exists.}
\label{fig:vtg_development_stresses}
\end{figure}

The distinction is necessary because a locally informative score can still induce a globally incomplete greedy policy. Conversely, a geometry that is not greedily navigable may still have substantial value if it sharply reduces a bounded frontier. We therefore propose the local-navigability curve
\[
N(k,B)=\bigl(\text{verified success},\text{route stretch},\text{states expanded},\text{verifier cost}\bigr),
\]
where $k$ is the amount of local branching/information permitted and $B$ is the resource ceiling used by the policy evaluation.

\subsection{Phase-0/1 falsification experiment}
The next mathematical epistemic cut is frozen as:
\begin{center}
\fbox{\parbox{0.91\linewidth}{\centering
\textbf{Next mathematical epistemic cut:} Does verifier-defined mathematical reachability admit a succinct, constructible, leakage-free and reusable directed reachability representation whose local queries reduce total verified proof-search cost on held-out theorem families?}}
\end{center}
The first experiment must not use an LLM as the solver. A finite Lean environment is split into: (a) an evaluator-only gold transition/proof universe; (b) a solver-visible partial map with hidden gold routes and intermediate labels removed; (c) a geometry-construction/training split; and (d) held-out theorem families. Every operational state is required to bind the future-relevant Lean environment/proof-state components needed for the registered transition semantics.

The registered controls are ordinary graph search/best-first, a gold-distance oracle available only to the evaluator, simple structural heuristics and a learned proof-progress baseline where reproducibly available. Candidate geometries are introduced one family at a time---for example directed/quasimetric or reachability representations, followed by multi-representation/quotient-aware forms only if simpler candidates leave a residual. Each candidate declares its constructibility class and nontriviality contract. Each geometry is tested under strict greedy, bounded-local branching and geometry-guided best-first policies. Geometry construction, coordinate/storage cost, local query cost, abstraction/refinement, independence checking and verifier cost are included in the resource account.

A semantic quotient used for acceleration must separately report whether it is an exact navigation quotient or an over-approximation requiring route lifting/refinement. A learned geometry must bind its behavior policy, training/support coverage, leakage checks and fresh subjects. Gold routes, gold distance labels and future map edges remain evaluator-only. The lane is declared unsupported in the registered scope if, after hidden-route leakage controls and total construction/storage/query/verification cost, no candidate geometry improves verified proof success or search cost over strong controls on held-out theorem families. Narrower outcomes are allowed, for example \path{LOCALLY_INFORMATIVE_NOT_GREEDILY_NAVIGABLE}, \path{ATLAS_NAVIGABILITY_SUPPORTED_GLOBAL_GEOMETRY_UNSUPPORTED}, \path{NONTRIVIAL_GEOMETRY_NOT_COST_EFFECTIVE}, or \path{NO_SAFE_NAVIGATION_QUOTIENT_AT_TESTED_SCALE}. Only after a positive Phase-1 signal would flow, diffusion, path-integral or conductance dynamics become justified experiments.

\paragraph{Scope boundary.}
The Phase-0/1 theory is deterministic. If later operators are stochastic or adversarial, the transition substrate must be upgraded to an appropriate stochastic kernel, MDP or game semantics and the reachability quantifier must be declared; a deterministic edge geometry must not silently absorb stochasticity. Likewise, on infinite state spaces an infimum path value need not be attained and a geodesic need not exist. The word \emph{geometry} is therefore used broadly; metric, quasimetric, geodesic, manifold or categorical-fibration properties are claimed only when their respective axioms are separately established.

\paragraph{Relationship to the broader Orion mechanics.}
Verified Transformation Geometry is the mathematical instantiation of the broader proposal-only Verified Solver Compilation layer. If navigation fails, Mechanic Differential Diagnosis asks whether the obstruction belongs to specification, state sufficiency, map coverage, navigation quotient/abstraction, representation/portal, scale, metric, dynamics, operator basis, verifier, interface, support/OOD, model/tool floor or compilation cost. A representation or method change is then requested by its desired transformation effect and must return through the same specification, preservation, amalgamation and proof-assurance gates described above. This preserves the design goal of aggressive search with conservative authority.